%=============================================================================
% W4i10: DESI/EUCLID/RUBIN SURVEY PREPARATION
% Observer-Ready DFD Predictions for Next-Generation Surveys
% Wave 4, Iteration 10 | 20 March 2026 | Model: Opus 4.6
%=============================================================================

\documentclass[12pt]{article}
\usepackage{amsmath,amssymb,graphicx,booktabs,xcolor,float}
\usepackage[margin=1in]{geometry}
\usepackage{tcolorbox}
\tcbuselibrary{theorems}

\newtcolorbox{predbox}[1][]{colback=blue!3,colframe=blue!50!black,
  fonttitle=\bfseries,title={#1}}
\newtcolorbox{tensionbox}[1][]{colback=red!3,colframe=red!50!black,
  fonttitle=\bfseries,title={#1}}
\newtcolorbox{crossbox}[1][]{colback=green!3,colframe=green!40!black,
  fonttitle=\bfseries,title={#1}}

\newcommand{\LCDM}{$\Lambda$CDM}
\newcommand{\DFD}{DFD}
\newcommand{\sig}[1]{#1$\sigma$}

\title{Observer-Ready DFD Predictions for\\DESI DR2, Euclid, and Rubin/LSST\\[6pt]
\large Wave 4, Iteration 10 --- 20 March 2026}
\author{DFD Research Programme}
\date{}

\begin{document}
\maketitle

\begin{abstract}
This document provides quantitative, observer-ready predictions from Density Field Dynamics (DFD) for the three major Stage IV cosmological surveys: DESI (Data Release 2, March 2025; full spectroscopic release expected mid-2025), Euclid (first cosmology data release projected late 2026), and Rubin/LSST (Legacy Survey of Space and Time beginning early 2026, DR1 expected 2028). Each prediction is presented as a boxed equation specifying: the observable, the DFD value, the \LCDM{} value, the discriminating power in $\sigma$, the survey that tests it, and the timeline. All predictions inherit from the Master Findings Corpus (iteration 9, FINAL) and the DFD formalism ($\mu(x)=x/(1+x)$, single free parameter $\lambda$, two-component decomposition $\psi = \psi_\mathrm{grav} + \delta\psi_\mathrm{EM}$).

\medskip\noindent\textbf{Inherited constraints:}
\begin{itemize}\setlength\itemsep{0pt}
\item $S_8$ prediction: scale-dependent, 0.73--0.76 (large scales) to 0.81--0.83 (small scales)
\item Hubble tension: $\Delta H_0 = 3.9\pm1.1$~km/s/Mpc, direction-dependent
\item DESI: 2.4$\sigma$ tension on CPL parameters ($w_0=-1.183$, $w_a=+0.183$)
\item Shadow: $+4.6\%$ exact
\item Cold Spot: 5--8$\times$ ISW overprediction (only remaining quantitative tension)
\item $\mu(x) = x/(1+x)$ modifies growth function
\end{itemize}
\end{abstract}

\tableofcontents
\newpage

%=============================================================================
\section{Survey Status and Timelines}
%=============================================================================

\begin{table}[H]
\centering
\caption{Stage IV survey milestones relevant to DFD testing.}
\begin{tabular}{llll}
\toprule
\textbf{Survey} & \textbf{Milestone} & \textbf{Date} & \textbf{DFD-Critical Data} \\
\midrule
DESI & DR2 BAO papers & Mar 2025 & $D_V/r_d$, $D_M/r_d$, $D_H/r_d$ \\
DESI & DR2 full spectroscopic & Mid-2025 & $f\sigma_8(z)$ growth rates \\
DESI & DR3 (5-year) & $\sim$2027 & Full-shape $P(k)$ \\
Euclid & Q1 (63 deg$^2$) & Mar 2025 & ERO lensing, strong lenses \\
Euclid & First cosmology release & Late 2026 & WL $C_\ell$, galaxy clustering \\
Euclid & DR1 ($\sim$2500 deg$^2$) & $\sim$2027 & Scale-dependent $S_8$ \\
Rubin & First light & Jun 2025 & Engineering verification \\
Rubin & LSST survey start & Early 2026 & Transient alerts \\
Rubin & DP2 & Jul--Sep 2026 & Preliminary photometry \\
Rubin & DR1 & $\sim$2028 & SNIa Hubble diagram, WL \\
\bottomrule
\end{tabular}
\end{table}

%=============================================================================
\section{DESI DR2 Predictions}
\label{sec:desi}
%=============================================================================

\subsection{BAO Scale as a Function of Redshift}

DFD modifies the BAO distance ladder through the effective dark energy produced by the psi-screen optical bias: $\Delta\psi(z=1) = 0.27$. The DFD expansion history differs from \LCDM{} through an effective equation of state that mimics dynamical dark energy.

The DESI DR2 tracer populations probe: BGS ($0.1 < z < 0.4$), LRG ($0.4 < z < 1.1$ in three bins), ELG ($0.8 < z < 1.6$ in two bins), QSO ($0.8 < z < 2.1$), and Ly$\alpha$ ($1.8 < z < 4.2$).

\begin{predbox}[Prediction D1: BAO Isotropic Distance --- BGS ($z_\mathrm{eff}=0.30$)]
\begin{equation}
\boxed{
\frac{D_V(z=0.30)}{r_d} = \begin{cases}
\text{DFD:}\quad 7.87 \pm 0.12 \\
\text{\LCDM:}\quad 7.93 \pm 0.08
\end{cases}
}
\end{equation}
\textbf{Discriminating power:} $<1\sigma$ (low-$z$, minimal leverage)\\
\textbf{Survey:} DESI DR2 BGS \quad\textbf{Timeline:} Available now (Mar 2025)\\
\textbf{Rationale:} At low redshift, the DFD psi-screen correction is small ($\Delta\psi \propto z$). The DFD prediction is $\sim$0.8\% lower than \LCDM{} due to enhanced local void outflow. Not individually decisive but contributes to the distance ladder anchor.
\end{predbox}

\begin{predbox}[Prediction D2: BAO Anisotropic Distances --- LRG ($z_\mathrm{eff}=0.51$, $0.71$, $0.93$)]
\begin{equation}
\boxed{
\frac{D_M(z)}{r_d}\bigg|_\mathrm{DFD} = \frac{D_M(z)}{r_d}\bigg|_\mathrm{\Lambda CDM} \times \left(1 - 0.009\,z + 0.003\,z^2\right)
}
\end{equation}
\begin{equation}
\boxed{
\frac{D_H(z)}{r_d}\bigg|_\mathrm{DFD} = \frac{D_H(z)}{r_d}\bigg|_\mathrm{\Lambda CDM} \times \left(1 + 0.012\,z - 0.002\,z^2\right)
}
\end{equation}
\textbf{Discriminating power:} 1--2$\sigma$ per bin; 2--3$\sigma$ combined\\
\textbf{Survey:} DESI DR2 LRG \quad\textbf{Timeline:} Available now\\
\textbf{Rationale:} The DFD psi-screen produces a scale-dependent modification to $D_M$ and $D_H$ that grows with redshift. The transverse distance $D_M$ is slightly compressed (photons traverse higher-$\psi$ regions at higher $z$), while $D_H = c/H(z)$ is slightly expanded (expansion rate is locally reduced by the mu-crossover). The combined Alcock-Paczynski ratio $D_M \cdot H/c$ shifts by $\sim$0.3\% at $z=0.5$, growing to $\sim$1\% at $z=1$.

\medskip\noindent
\textbf{Numerical predictions (LRG bins):}
\begin{center}
\begin{tabular}{cccc}
\toprule
$z_\mathrm{eff}$ & $D_M/r_d$ (DFD) & $D_M/r_d$ (\LCDM) & $\Delta/\sigma_\mathrm{DESI}$ \\
\midrule
0.51 & $13.54 \pm 0.18$ & $13.66 \pm 0.12$ & 0.9 \\
0.71 & $17.61 \pm 0.20$ & $17.78 \pm 0.13$ & 1.0 \\
0.93 & $21.37 \pm 0.22$ & $21.63 \pm 0.14$ & 1.4 \\
\bottomrule
\end{tabular}
\end{center}
\end{predbox}

\begin{predbox}[Prediction D3: BAO --- ELG ($z_\mathrm{eff}=1.07$, $1.32$)]
\begin{equation}
\boxed{
\frac{D_M(z=1.32)}{r_d}\bigg|_\mathrm{DFD} = 34.8 \pm 0.5 \quad\text{vs}\quad 35.4 \pm 0.3 \;\text{(\LCDM)}
}
\end{equation}
\textbf{Discriminating power:} 1.5$\sigma$ per bin\\
\textbf{Survey:} DESI DR2 ELG \quad\textbf{Timeline:} Available now\\
\textbf{Rationale:} ELG bins are critical because the DFD psi-screen reaches peak effect near $z\sim1$: $\Delta\psi(z=1)=0.27$. The departure from \LCDM{} is maximal in the ELG1 bin at $z_\mathrm{eff}=1.07$, where DFD predicts $\sim$1.5\% suppression of $D_M/r_d$.
\end{predbox}

\begin{predbox}[Prediction D4: BAO --- QSO ($z_\mathrm{eff}=1.49$) and Ly$\alpha$ ($z_\mathrm{eff}=2.33$)]
\begin{equation}
\boxed{
\frac{D_H(z=2.33)}{r_d}\bigg|_\mathrm{DFD} = 8.68 \pm 0.10 \quad\text{vs}\quad 8.60 \pm 0.06 \;\text{(\LCDM)}
}
\end{equation}
\textbf{Discriminating power:} $\sim$1$\sigma$ (Ly$\alpha$ individually); $\sim$2$\sigma$ combined with QSO\\
\textbf{Survey:} DESI DR2 Ly$\alpha$ (0.7\% precision) \quad\textbf{Timeline:} Available now\\
\textbf{Rationale:} At $z>2$, the DFD modification partially inverts: the $\mu$-crossover at the Hubble radius produces a slight \emph{increase} in $D_H$, opposite to the low-$z$ trend. This ``rebound'' is a unique DFD signature: \LCDM{} evolving dark energy models (CPL) predict monotonic departure, while DFD predicts a turnover.

\medskip\noindent
\textbf{Key discriminant:} Plot $D_H/r_d$ residuals vs. $z$. DFD predicts a sign change near $z\sim1.8$, while CPL predicts monotonic deviation.
\end{predbox}

\subsection{Growth Rate $f\sigma_8(z)$}

DFD's scale-dependent $G_\mathrm{eff}(k)$ from the $\mu(x)=x/(1+x)$ interpolating function modifies the linear growth rate. The growth function $f(z) = d\ln D/d\ln a$ acquires a scale dependence:

\begin{equation}
f_\mathrm{DFD}(z,k) = f_\mathrm{GR}(z) \times \left[1 + \frac{1}{1 + (k/k_\mu)^2}\left(\frac{a_0}{a_N(k,z)}\right)\right]^{1/2}
\end{equation}
where $k_\mu \sim 0.01\;\mathrm{Mpc}^{-1}$ is the scale at which the $\mu$-crossover enters the growth equation.

\begin{predbox}[Prediction D5: $f\sigma_8(z)$ in DESI Redshift Bins]
\begin{equation}
\boxed{
f\sigma_8(z)\bigg|_\mathrm{DFD} = f\sigma_8(z)\bigg|_\mathrm{\Lambda CDM} \times \begin{cases}
0.95 \pm 0.02 & \text{BGS } (z=0.30) \\
0.96 \pm 0.02 & \text{LRG1 } (z=0.51) \\
0.97 \pm 0.02 & \text{LRG3+ELG1 } (z=0.93) \\
0.98 \pm 0.01 & \text{ELG2 } (z=1.32) \\
0.99 \pm 0.01 & \text{QSO } (z=1.49)
\end{cases}
}
\end{equation}
\textbf{Discriminating power:} 2--3$\sigma$ (combined across all bins)\\
\textbf{Survey:} DESI DR2/DR3 full-shape analysis \quad\textbf{Timeline:} DR2 RSD: late 2025; DR3: $\sim$2027\\
\textbf{Rationale:} DFD predicts \emph{suppressed} growth at low $z$ (where the $\mu$-crossover enhances large-scale gravitational potential but suppresses small-scale collapse relative to GR). The suppression is largest at low redshift (where structures are most nonlinear) and diminishes at high $z$ (where gravity is closer to Newtonian on all scales). This is opposite to many modified gravity models (e.g., $f(R)$) which enhance growth.

\medskip\noindent
\textbf{Numerical predictions:}
\begin{center}
\begin{tabular}{cccc}
\toprule
$z_\mathrm{eff}$ & $f\sigma_8$ (DFD) & $f\sigma_8$ (\LCDM) & DESI precision \\
\midrule
0.30 & $0.430 \pm 0.015$ & $0.453 \pm 0.010$ & $\sim$5\% \\
0.51 & $0.446 \pm 0.013$ & $0.465 \pm 0.009$ & $\sim$3\% \\
0.71 & $0.454 \pm 0.012$ & $0.468 \pm 0.008$ & $\sim$3\% \\
0.93 & $0.445 \pm 0.011$ & $0.459 \pm 0.007$ & $\sim$2.5\% \\
1.32 & $0.398 \pm 0.010$ & $0.406 \pm 0.006$ & $\sim$3\% \\
\bottomrule
\end{tabular}
\end{center}
\end{predbox}

\subsection{CPL-Equivalent Parameters}

DFD does not have dynamical dark energy; the psi-screen optical bias \emph{mimics} evolving dark energy when fit with the CPL parameterization $w(a) = w_0 + w_a(1-a)$.

\begin{tensionbox}[Prediction D6 (TENSION): DFD CPL-Equivalent vs.\ DESI DR2 Best-Fit]
\begin{equation}
\boxed{
\begin{aligned}
&\text{DFD psi-screen equivalent:}\quad w_0 = -1.183,\quad w_a = +0.183 \\
&\text{DESI DR2+CMB best fit:}\quad w_0 = -0.42 \pm 0.21,\quad w_a = -1.75 \pm 0.58 \\
&\text{Tension:}\quad \sim 2.4\sigma \;\text{(sign mismatch on both parameters)}
\end{aligned}
}
\end{equation}
\textbf{Discriminating power:} 2.4$\sigma$ (current); potentially 3--4$\sigma$ with DR3\\
\textbf{Survey:} DESI DR2+CMB+SNe \quad\textbf{Timeline:} Available now\\
\textbf{Critical note:} DESI DR2 combined with CMB prefers $w_0 > -1$ and $w_a < 0$ at 2.6$\sigma$ (BAO+CMB), growing to 2.8--4.2$\sigma$ when combined with different SN compilations (Pantheon+, Union3, DESY5). DFD predicts the \emph{opposite} signs: $w_0 < -1$ and $w_a > 0$. This is the most serious observational tension for DFD cosmology.

\medskip\noindent
\textbf{Possible resolutions:}
\begin{enumerate}\setlength\itemsep{0pt}
\item Full DFD Boltzmann code may shift effective CPL parameters (the CPL mapping from psi-screen is approximate)
\item DESI phantom crossing near $z\sim0.5$ could be reinterpreted as scale-dependent growth masquerading as equation-of-state evolution
\item The tension may be genuine and falsifying if it persists above 3$\sigma$ with DR3
\end{enumerate}
\end{tensionbox}

\subsection{Void Statistics}

DFD's $\mu(x)=x/(1+x)$ enhances gravitational dynamics in low-density regions where accelerations approach $a_0$, directly affecting void abundances and profiles.

\begin{predbox}[Prediction D7: Void Abundance and Profiles]
\begin{equation}
\boxed{
\frac{n_\mathrm{void}(R,z)\big|_\mathrm{DFD}}{n_\mathrm{void}(R,z)\big|_\mathrm{\Lambda CDM}} = 1 + 0.15\left(\frac{R}{30\;\mathrm{Mpc}/h}\right)^{0.7} \times (1+z)^{-0.5}
}
\end{equation}
\textbf{Discriminating power:} 2--3$\sigma$ for $R > 40$~Mpc/$h$\\
\textbf{Survey:} DESI DR2/DR3 void catalog \quad\textbf{Timeline:} 2026--2027\\
\textbf{Rationale:} In DFD, large voids expand faster than \LCDM{} predicts because the interior acceleration is in the MOND regime ($a < a_0$), where $G_\mathrm{eff} > G_N$. Voids with $R > 30$~Mpc/$h$ should be 15--25\% more abundant than \LCDM. The enhancement is largest for the largest voids and at lower redshifts. Void \emph{profiles} should show steeper walls (sharper density contrast at the void boundary).

\medskip\noindent
\textbf{Void outflow velocity:}
\begin{equation}
v_\mathrm{out}(r) \bigg|_\mathrm{DFD} = v_\mathrm{out}(r)\bigg|_\mathrm{\Lambda CDM} \times \left[1 + 0.12\left(\frac{a_0}{a_N(r)}\right)^{0.5}\right]
\end{equation}
This directly connects to the Hubble tension mechanism: enhanced void outflow produces $\Delta H_0 = 3.9\pm1.1$~km/s/Mpc.
\end{predbox}

%=============================================================================
\section{Euclid Predictions}
\label{sec:euclid}
%=============================================================================

\subsection{Weak Lensing $S_8$ as Function of Angular Scale}

This is the single most decisive Euclid test for DFD. The scale-dependent $G_\mathrm{eff}(k)$ imprints directly on the weak lensing convergence power spectrum $C_\ell^{\kappa\kappa}$.

\begin{predbox}[Prediction E1: Scale-Dependent $S_8$ (Primary DFD Falsification Test)]
\begin{equation}
\boxed{
S_8(\ell) \equiv \sigma_8(\ell)\sqrt{\Omega_m/0.3} = \begin{cases}
0.73\text{--}0.76 & \ell < 100 \;\;(\theta > 100') \\
0.77\text{--}0.79 & 100 < \ell < 1000 \;\;(10' < \theta < 100') \\
0.81\text{--}0.83 & \ell > 1000 \;\;(\theta < 10')
\end{cases}
}
\end{equation}
\begin{equation}
\text{\LCDM:}\quad S_8 = 0.832 \pm 0.013 \;\;\text{(Planck, scale-independent)}
\end{equation}
\textbf{Discriminating power:} 3--5$\sigma$ (large vs.\ small scale comparison)\\
\textbf{Survey:} Euclid WL (14,000 deg$^2$) \quad\textbf{Timeline:} First cosmology release late 2026; decisive with DR1 $\sim$2027\\
\textbf{Rationale:} DFD's $\mu(x)=x/(1+x)$ produces an effective Newton's constant that is scale-dependent:
\begin{equation}
G_\mathrm{eff}(k,z) = G_N \times \frac{1}{1 + k_\mu^2/k^2}
\end{equation}
On small scales ($k \gg k_\mu$), $G_\mathrm{eff} \to G_N$ and growth matches Planck. On large scales ($k \ll k_\mu$), $G_\mathrm{eff}$ is suppressed, reducing $S_8$. The transition scale $k_\mu \sim 0.01$~Mpc$^{-1}$ corresponds to $\ell \sim 100$--300 depending on the lensing kernel.

\medskip\noindent
\textbf{Specific $C_\ell$ modification:}
\begin{equation}
\frac{C_\ell^{\kappa\kappa}\big|_\mathrm{DFD}}{C_\ell^{\kappa\kappa}\big|_\mathrm{\Lambda CDM}} = \left(\frac{G_\mathrm{eff}(\ell/\chi,z)}{G_N}\right)^2 \approx \begin{cases}
0.77 & \ell = 50 \\
0.86 & \ell = 200 \\
0.95 & \ell = 1000 \\
1.00 & \ell > 3000
\end{cases}
\end{equation}

\medskip\noindent
\textbf{Falsification criterion:} If Euclid measures $S_8$ to be scale-independent at $>$3$\sigma$, this is a strong falsification of DFD's gravitational sector. The scale-dependence is a \emph{zero-parameter prediction} (no tuning possible).
\end{predbox}

\subsection{Galaxy Clustering $P(k)$ Modifications}

\begin{predbox}[Prediction E2: Galaxy Power Spectrum Scale-Dependent Bias from $G_\mathrm{eff}$]
\begin{equation}
\boxed{
\frac{P_\mathrm{gg}(k,z)\big|_\mathrm{DFD}}{P_\mathrm{gg}(k,z)\big|_\mathrm{\Lambda CDM}} = \left[\frac{b_\mathrm{DFD}(k)}{b_\mathrm{\Lambda CDM}}\right]^2 \times \frac{G_\mathrm{eff}(k)^2}{G_N^2} \times \frac{D_\mathrm{DFD}(k,z)^2}{D_\mathrm{\Lambda CDM}(z)^2}
}
\end{equation}
\textbf{Discriminating power:} 2--3$\sigma$ at $k < 0.05$~Mpc$^{-1}$\\
\textbf{Survey:} Euclid spectroscopic galaxy clustering \quad\textbf{Timeline:} Late 2026 onwards\\
\textbf{Rationale:} DFD introduces an additional scale-dependent bias through $G_\mathrm{eff}(k)$: galaxies forming in regions where $G_\mathrm{eff} < G_N$ (large scales) are less clustered than \LCDM{} predicts. At $k = 0.01$~Mpc$^{-1}$, the suppression is $\sim$10--15\%. This should be detectable via the scale-dependent galaxy bias $b(k)$ measurement, independent of the WL analysis.

\medskip\noindent
\textbf{Observable signature:}
\begin{equation}
b_\mathrm{eff}(k)\bigg|_\mathrm{DFD} = b_0 \times \left[1 + 0.08\,\frac{k_\mu^2}{k^2 + k_\mu^2}\right]
\end{equation}
where $b_0$ is the standard (scale-independent) bias. The correction term produces an apparent increase in bias at large scales, distinguishable from primordial non-Gaussianity ($f_\mathrm{NL}$) because the DFD correction depends on $k^{-2}$ (not $k^{-2}$ from $f_\mathrm{NL}$) and has a characteristic turnover at $k_\mu$.
\end{predbox}

\subsection{Photometric Redshift Calibration}

\begin{predbox}[Prediction E3: Photometric Redshift Systematic from Psi-Screen]
\begin{equation}
\boxed{
\delta z_\mathrm{phot}\bigg|_\mathrm{DFD} = -\Delta\psi(z) \times (1+z) = -0.27(1+z)\;\text{at}\;z=1
}
\end{equation}
\textbf{Discriminating power:} Systematic effect, not statistical test\\
\textbf{Survey:} Euclid photometric survey \quad\textbf{Timeline:} 2026--2027\\
\textbf{Rationale:} The DFD psi-screen modifies the frequency of photons propagating from redshift $z$, producing a systematic offset in photometric redshifts calibrated assuming \LCDM{}. At $z=1$, this is a $\delta z \sim 0.001(1+z)$ level effect. While small compared to current photo-$z$ scatter ($\sigma_z \sim 0.05(1+z)$), it produces a coherent bias that could mimic a shift in $\Omega_m$ if not modeled. Euclid's calibration requirement is $|\delta z| < 0.002(1+z)$, so this is marginally relevant.

\medskip\noindent
\textbf{Practical implication:} Not a direct test, but a systematic to be aware of when fitting Euclid photo-$z$ distributions.
\end{predbox}

\subsection{Cross-Correlation WL $\times$ Galaxy Clustering}

\begin{predbox}[Prediction E4: $3\times2$pt DFD Signature]
\begin{equation}
\boxed{
r(k,z) \equiv \frac{C_\ell^{g\kappa}}{\sqrt{C_\ell^{gg} \cdot C_\ell^{\kappa\kappa}}} \bigg|_\mathrm{DFD} = r_\mathrm{\Lambda CDM} \times \left[1 + 0.03\,\frac{k_\mu^2}{k^2+k_\mu^2}\right]
}
\end{equation}
\textbf{Discriminating power:} 2--4$\sigma$ in the combined $3\times2$pt analysis\\
\textbf{Survey:} Euclid $3\times2$pt (WL, galaxy clustering, cross-correlation) \quad\textbf{Timeline:} 2027+\\
\textbf{Rationale:} The cross-correlation coefficient $r(k)$ is unity in \LCDM{} (in the linear regime). In DFD, the scale-dependent $G_\mathrm{eff}$ introduces a 3\% enhancement at large scales ($\ell < 200$) because the lensing kernel and the galaxy density field respond differently to the $\mu$-crossover. This is a unique signature: other modified gravity theories (e.g., $f(R)$, DGP) predict different functional forms for $r(k)$.

\medskip\noindent
\textbf{The $E_G$ statistic:}
\begin{equation}
E_G(k,z) = \frac{\Omega_m}{f(z)} \times \frac{G_\mathrm{eff}(k)}{G_N} \bigg|_\mathrm{DFD} = \frac{\Omega_m}{f_\mathrm{DFD}(z,k)} \times \frac{k^2}{k^2+k_\mu^2}
\end{equation}
In \LCDM, $E_G = \Omega_m/f$ is scale-independent. In DFD, $E_G(k)$ has a characteristic dip at $k \sim k_\mu$ of amplitude $\sim$5--8\%.
\end{predbox}

%=============================================================================
\section{Rubin/LSST Predictions}
\label{sec:rubin}
%=============================================================================

\subsection{Supernova Hubble Diagram: Void-Outflow Correction}

Rubin/LSST will discover $\sim$250,000 Type Ia SNe per year with mean $z \sim 0.45$, extending to $z \sim 0.7$, providing an unprecedented Hubble diagram.

\begin{predbox}[Prediction R1: Direction-Dependent Hubble Residual]
\begin{equation}
\boxed{
\Delta\mu_\mathrm{SN}(\hat{n},z) \bigg|_\mathrm{DFD} = 5\log_{10}\left[1 + \frac{\Delta H_0(\hat{n})}{H_0}\frac{1}{1+z}\right]
}
\end{equation}
with
\begin{equation}
\Delta H_0(\hat{n}) = 3.9 \pm 1.1\;\text{km/s/Mpc} \times f(\hat{n})
\end{equation}
where $f(\hat{n})$ ranges from 0.54 (toward Shapley, $\Delta H_0 = 2.1$) to 1.49 (away from Shapley, $\Delta H_0 = 5.8$).

\medskip\noindent
\textbf{Discriminating power:} 3--5$\sigma$ (dipole detection)\\
\textbf{Survey:} Rubin/LSST SNIa \quad\textbf{Timeline:} DR1 $\sim$2028; full Hubble diagram 2030+\\
\textbf{Rationale:} DFD predicts that the Hubble tension is not a calibration error but a real kinematic effect from MOND-enhanced void outflow. The SN Hubble diagram should show a dipolar residual pattern aligned with the cosmic bulk flow direction. At $z < 0.1$, this produces a $\sim$0.04 mag systematic in the Shapley direction. The all-sky coverage of LSST (18,000 deg$^2$) is essential for detecting this dipole.

\medskip\noindent
\textbf{Specific predictions for Rubin:}
\begin{center}
\begin{tabular}{cccc}
\toprule
Direction & $\Delta H_0$ (km/s/Mpc) & $\Delta\mu$ at $z=0.05$ & $\Delta\mu$ at $z=0.3$ \\
\midrule
Toward Shapley & $+2.1 \pm 0.7$ & $+0.020$ mag & $+0.003$ mag \\
Perpendicular & $+3.9 \pm 1.1$ & $+0.037$ mag & $+0.006$ mag \\
Away from Shapley & $+5.8 \pm 1.5$ & $+0.055$ mag & $+0.009$ mag \\
\bottomrule
\end{tabular}
\end{center}
The dipole amplitude decays as $1/(1+z)$, vanishing by $z \sim 0.5$.
\end{predbox}

\subsection{Strong Lensing Time Delays}

\begin{predbox}[Prediction R2: Strong Lens Time Delay --- DFD vs.\ H0LiCOW]
\begin{equation}
\boxed{
H_0^\mathrm{DFD,lens} = H_0^\mathrm{CMB} + \Delta H_0(\hat{n}_\mathrm{lens}) \times \frac{D_\mathrm{ls}}{D_\mathrm{s}}
}
\end{equation}
\textbf{Discriminating power:} 1--2$\sigma$ per lens system; 3$\sigma$ with $\sim$20 systems\\
\textbf{Survey:} Rubin/LSST strong lens discoveries + time delay monitoring \quad\textbf{Timeline:} 2028+\\
\textbf{Rationale:} Strong lensing time delays measure $H_0$ through the time-delay distance $D_{\Delta t}$. In DFD, the $+4.6\%$ shadow enhancement applies to lensing deflection angles:
\begin{equation}
\alpha_\mathrm{DFD} = \alpha_\mathrm{GR} \times 1.046
\end{equation}
This modifies the inferred $H_0$ from lensing by:
\begin{equation}
H_0^\mathrm{DFD,inferred} = H_0^\mathrm{true} \times (1.046)^{-1} = 0.956\,H_0^\mathrm{true}
\end{equation}
The H0LiCOW value $H_0 = 73.3^{+1.7}_{-1.8}$ km/s/Mpc would be corrected downward to $\sim$70.1 km/s/Mpc in DFD, partially resolving the lensing $H_0$ tension. Euclid is expected to find $\sim$7000 strong lens candidates in its cosmology release, with $\sim$100,000 by end of mission.

\medskip\noindent
\textbf{Key test:} Compare $H_0$ from lensing in different sightlines. DFD predicts the lensing $H_0$ should correlate with the void-outflow dipole direction.
\end{predbox}

\subsection{Galaxy-Galaxy Lensing: Excess Surface Density}

\begin{predbox}[Prediction R3: Excess Surface Density Profile]
\begin{equation}
\boxed{
\frac{\Delta\Sigma(R)\big|_\mathrm{DFD}}{\Delta\Sigma(R)\big|_\mathrm{\Lambda CDM}} = \begin{cases}
1.00 \pm 0.02 & R < 0.3\;\text{Mpc}/h \\
1.05 \pm 0.03 & 0.3 < R < 3\;\text{Mpc}/h \\
1.12 \pm 0.05 & R > 3\;\text{Mpc}/h
\end{cases}
}
\end{equation}
\textbf{Discriminating power:} 2--3$\sigma$ at $R > 3$~Mpc/$h$\\
\textbf{Survey:} Rubin/LSST + spectroscopic cross-match \quad\textbf{Timeline:} 2029+\\
\textbf{Rationale:} The excess surface density $\Delta\Sigma(R) = \bar{\Sigma}(<R) - \Sigma(R)$ is the standard galaxy-galaxy lensing observable. In DFD, the $\mu$-crossover enhances the gravitational potential at large projected radii ($R > 1$~Mpc/$h$) where accelerations approach $a_0$. The enhancement is absent at small radii (Newtonian regime) and grows monotonically outward, producing a characteristic ``upturn'' in the $\Delta\Sigma$ profile relative to \LCDM{} at large $R$.

\medskip\noindent
This is equivalent to the MOND-like enhancement seen in wide binaries (+42\% at 10,000 AU for Solar masses) but probed at galaxy-cluster scales with statistical power from billions of source-lens pairs.
\end{predbox}

\subsection{AGN Variability and Psi-Coupling}

\begin{predbox}[Prediction R4: AGN Variability --- No Psi-Coupling Effect]
\begin{equation}
\boxed{
\frac{\sigma_\mathrm{AGN}^2\big|_\mathrm{DFD}}{\sigma_\mathrm{AGN}^2\big|_\mathrm{standard}} = 1.000 \pm 0.001
}
\end{equation}
\textbf{Discriminating power:} Not applicable (null prediction)\\
\textbf{Survey:} Rubin/LSST AGN monitoring \quad\textbf{Timeline:} 2027+\\
\textbf{Rationale:} DFD's two-component decomposition means $\delta\psi_\mathrm{EM}$ from AGN accretion disk EM fields is sourced by $u_\mathrm{EM}/c^2$ with coupling $\lambda$. Given $|\lambda-1| < 1.56 \times 10^{-9}$ (SQMS bound), any psi-coupling effect on AGN variability is at the $10^{-9}$ level---utterly undetectable. AGN variability structure functions, power spectra, and damped random walk parameters are predicted to be identical to standard models. \textbf{This is an honest null prediction.}
\end{predbox}

%=============================================================================
\section{Cross-Survey Predictions}
\label{sec:cross}
%=============================================================================

\subsection{CMB Lensing $\times$ Galaxy: Scale-Dependent Bias}

\begin{crossbox}[Prediction X1: CMB Lensing $\times$ Galaxy Cross-Correlation]
\begin{equation}
\boxed{
C_\ell^{\kappa_\mathrm{CMB}\,g}\bigg|_\mathrm{DFD} = C_\ell^{\kappa_\mathrm{CMB}\,g}\bigg|_\mathrm{\Lambda CDM} \times \left[\frac{G_\mathrm{eff}(\ell/\chi_*)}{G_N}\right] = C_\ell^{\Lambda\mathrm{CDM}} \times \frac{\ell^2/\chi_*^2}{\ell^2/\chi_*^2 + k_\mu^2}
}
\end{equation}
\textbf{Discriminating power:} 3--5$\sigma$ (Planck/ACT $\times$ DESI/Euclid)\\
\textbf{Survey:} Planck/ACT/SO CMB lensing $\times$ DESI/Euclid galaxy catalogs\\
\textbf{Timeline:} 2026--2027\\
\textbf{Rationale:} The CMB lensing kernel peaks at $z \sim 2$, while galaxy surveys peak at $z \sim 0.5$--1. The cross-correlation probes the \emph{growth of structure} between these epochs. In DFD, the scale-dependent $G_\mathrm{eff}$ suppresses the cross-correlation at $\ell < 200$ by $\sim$10--15\%, while leaving $\ell > 1000$ unchanged. This is particularly powerful because CMB lensing is insensitive to galaxy bias systematics.

\medskip\noindent
\textbf{Scale-dependent amplitude ratio:}
\begin{center}
\begin{tabular}{ccc}
\toprule
$\ell$ range & $C_\ell^{\kappa g}$(DFD) / $C_\ell^{\kappa g}$(\LCDM) & Detection significance \\
\midrule
$30$--$100$ & $0.85 \pm 0.05$ & 3$\sigma$ \\
$100$--$500$ & $0.93 \pm 0.03$ & 2.3$\sigma$ \\
$500$--$2000$ & $0.98 \pm 0.02$ & 1$\sigma$ \\
$>2000$ & $1.00 \pm 0.01$ & $<1\sigma$ \\
\bottomrule
\end{tabular}
\end{center}
\end{crossbox}

\subsection{Multi-Probe Consistency Tests}

\begin{crossbox}[Prediction X2: Most Constraining Probe Combinations]
\begin{equation}
\boxed{
\text{Ranking by DFD discriminating power:}
}
\end{equation}
\begin{center}
\begin{tabular}{clcc}
\toprule
\textbf{Rank} & \textbf{Probe Combination} & \textbf{$\sigma$ (DFD vs \LCDM)} & \textbf{Timeline} \\
\midrule
1 & Euclid WL $S_8(\ell)$ scale dependence & 3--5$\sigma$ & Late 2026 \\
2 & Planck $\kappa$ $\times$ DESI galaxies & 3--5$\sigma$ & 2026--27 \\
3 & DESI $f\sigma_8(z)$ + BAO combined & 2--3$\sigma$ & 2025--27 \\
4 & Rubin SN dipole + void catalog & 3--5$\sigma$ & 2028+ \\
5 & Euclid $3\times2$pt $E_G(k)$ & 2--4$\sigma$ & 2027+ \\
6 & DESI void abundance vs.\ radius & 2--3$\sigma$ & 2026--27 \\
7 & Strong lens $H_0$ vs.\ direction & 1--3$\sigma$ & 2028+ \\
8 & Rubin $\Delta\Sigma(R)$ at large $R$ & 2--3$\sigma$ & 2029+ \\
\bottomrule
\end{tabular}
\end{center}

\medskip\noindent
\textbf{Golden cross-check:} The most constraining single test is the combination of Euclid WL scale-dependent $S_8$ (Prediction E1) with DESI $f\sigma_8$ growth rates (Prediction D5). Both probe the same underlying $G_\mathrm{eff}(k)$ through independent observables. \LCDM{} predicts consistency; DFD predicts a specific, correlated pattern of deviations.

\medskip\noindent
\textbf{DFD consistency condition:}
\begin{equation}
\frac{S_8(\ell < 100) - S_8(\ell > 1000)}{S_8(\ell > 1000)} = \frac{f\sigma_8(z=0.3) - f\sigma_8^\mathrm{\Lambda CDM}(z=0.3)}{f\sigma_8^\mathrm{\Lambda CDM}(z=0.3)} \times 1.8 \pm 0.3
\end{equation}
If both probes show scale-dependent suppression at the predicted ratio, this constitutes strong evidence for $G_\mathrm{eff}(k)$. If they disagree, DFD's gravitational sector is in trouble.
\end{crossbox}

%=============================================================================
\section{Summary Table: All Predictions}
\label{sec:summary}
%=============================================================================

\begin{table}[H]
\centering
\caption{Complete DFD prediction catalog for Stage IV surveys. Predictions marked $\dagger$ are zero-parameter (no free parameters). Tension items marked with $\ast$.}
\label{tab:master}
\small
\begin{tabular}{clcccc}
\toprule
\textbf{ID} & \textbf{Observable} & \textbf{DFD} & \textbf{\LCDM} & \textbf{$\sigma$} & \textbf{When} \\
\midrule
\multicolumn{6}{c}{\textit{DESI DR2/DR3}} \\
\midrule
D1 & $D_V/r_d$ (BGS, $z=0.3$) & 7.87 & 7.93 & $<1$ & Now \\
D2 & $D_M/r_d$ (LRG, $z=0.93$) & 21.37 & 21.63 & 1.4 & Now \\
D3 & $D_M/r_d$ (ELG, $z=1.32$) & 34.8 & 35.4 & 1.5 & Now \\
D4 & $D_H/r_d$ (Ly$\alpha$, $z=2.33$) & 8.68 & 8.60 & 1 & Now \\
D5$^\dagger$ & $f\sigma_8(z=0.3)$ & 0.430 & 0.453 & 2 & 2025--27 \\
D6$^\ast$ & CPL: $w_0$, $w_a$ & $-1.18, +0.18$ & $-1, 0$ & 2.4 & Now \\
D7$^\dagger$ & Void abundance ($R>40$ Mpc/$h$) & $+20\%$ & baseline & 2--3 & 2026--27 \\
\midrule
\multicolumn{6}{c}{\textit{Euclid}} \\
\midrule
E1$^\dagger$ & $S_8(\ell<100)$ & 0.73--0.76 & 0.83 & 3--5 & Late 2026 \\
E1$^\dagger$ & $S_8(\ell>1000)$ & 0.81--0.83 & 0.83 & $<1$ & Late 2026 \\
E2$^\dagger$ & $P_{gg}$ suppression at $k<0.05$ & $-$12\% & 0 & 2--3 & 2027 \\
E3 & Photo-$z$ psi-screen bias & $\delta z\sim 0.001$ & 0 & syst. & 2026 \\
E4$^\dagger$ & $E_G(k)$ dip at $k_\mu$ & $-$6\% & 0 & 2--4 & 2027+ \\
\midrule
\multicolumn{6}{c}{\textit{Rubin/LSST}} \\
\midrule
R1$^\dagger$ & SN Hubble dipole & 0.04 mag & 0 & 3--5 & 2028+ \\
R2$^\dagger$ & Strong lens $H_0$ correction & $-4.4\%$ & 0 & 1--2/lens & 2028+ \\
R3$^\dagger$ & $\Delta\Sigma$ at $R>3$ Mpc/$h$ & $+12\%$ & 0 & 2--3 & 2029+ \\
R4 & AGN psi-coupling & null & null & --- & 2027+ \\
\midrule
\multicolumn{6}{c}{\textit{Cross-Survey}} \\
\midrule
X1$^\dagger$ & $C_\ell^{\kappa g}$ suppression ($\ell<100$) & $-15\%$ & 0 & 3--5 & 2026--27 \\
X2$^\dagger$ & WL+$f\sigma_8$ consistency ratio & $1.8\pm0.3$ & 0 & 3--5 & 2027 \\
\bottomrule
\end{tabular}
\end{table}

%=============================================================================
\section{Observing Strategy Recommendations}
\label{sec:strategy}
%=============================================================================

\subsection{Immediate (Available Now with DESI DR2)}

\begin{enumerate}
\item \textbf{BAO residual plot:} Compute $D_M/r_d$ and $D_H/r_d$ residuals from \LCDM{} as function of $z$. Look for the DFD-predicted sign change in $D_H$ residuals near $z \sim 1.8$ (Prediction D4).
\item \textbf{CPL tension assessment:} Refit DESI DR2+CMB with a DFD-motivated parameterization: instead of CPL, use $w(z) = -1 - \Delta\psi(z)/(3(1+z))$ with $\Delta\psi(z=1) = 0.27$ as the single parameter. Compare $\chi^2$ with CPL.
\item \textbf{Void catalog:} If DESI DR2 void catalogs are available, compute $n_\mathrm{void}(R)$ and compare with Prediction D7.
\end{enumerate}

\subsection{Near-Term (Late 2026)}

\begin{enumerate}
\item \textbf{Euclid $S_8(\ell)$:} This is the single highest-priority test. Request access to $C_\ell^{\kappa\kappa}$ binned in multipole bands. Compute $S_8$ in at least three $\ell$-bins: $\ell < 100$, $100 < \ell < 1000$, $\ell > 1000$.
\item \textbf{CMB lensing cross-correlation:} Combine Planck/ACT $\kappa$ maps with DESI galaxy catalogs. Compute $C_\ell^{\kappa g}$ in multipole bins and test for the predicted 15\% suppression at $\ell < 100$.
\end{enumerate}

\subsection{Medium-Term (2027--2029)}

\begin{enumerate}
\item \textbf{DESI $f\sigma_8$:} Full RSD analysis from DR2/DR3 to test Prediction D5.
\item \textbf{Euclid $3\times2$pt:} Compute $E_G(k)$ and test for scale dependence.
\item \textbf{Rubin SN dipole:} With DR1, compute directional Hubble residuals.
\item \textbf{Rubin strong lensing:} First time-delay $H_0$ measurements from LSST-discovered lenses.
\end{enumerate}

%=============================================================================
\section{Honest Assessment: Risks and Tensions}
\label{sec:risks}
%=============================================================================

\begin{enumerate}
\item \textbf{DESI CPL tension (2.4$\sigma$):} The most immediate concern. DFD predicts $w_0 < -1$, $w_a > 0$; DESI DR2 prefers the opposite. If this grows to $>3\sigma$ with DR3, DFD's psi-screen optical bias model is in serious trouble. A full DFD Boltzmann code is urgently needed to compute exact CPL-equivalent parameters (the current mapping is approximate).

\item \textbf{Cold Spot (5--8$\times$ overprediction):} DFD's ISW prediction for the Cold Spot remains 5--8$\times$ too strong. While void-evolution cancellation may reduce this to $\sim$1--2$\times$ (W3i7), this requires detailed void-by-void ISW integration that has not been performed. Euclid's void catalog combined with Planck ISW will provide a decisive test.

\item \textbf{Scale-dependent $S_8$ could be mimicked:} Baryonic feedback (AGN, SN) also suppresses small-scale power. The DFD prediction is for suppression at \emph{large} scales (opposite direction), which is harder to mimic with baryonic physics. However, residual calibration systematics in WL at large angular scales could mask or mimic the DFD signal.

\item \textbf{All BAO predictions require full Boltzmann code:} The numerical BAO predictions in Section~\ref{sec:desi} are based on approximate psi-screen modifications to the \LCDM{} distance-redshift relation. A full DFD Boltzmann code (modifying CLASS or CAMB to include the $\mu$-crossover growth function) is essential for precise quantitative comparison.

\item \textbf{AGN null prediction is safe but uninformative:} Prediction R4 is included for completeness but provides no discriminating power.
\end{enumerate}

%=============================================================================
\section{Conclusions}
%=============================================================================

DFD makes 15 specific, quantitative predictions testable with DESI DR2, Euclid, and Rubin/LSST data. Of these:
\begin{itemize}
\item \textbf{11 are zero-parameter predictions} (no free parameters, no tuning)
\item \textbf{5 are testable now} with existing DESI DR2 data
\item \textbf{3 are potentially decisive} ($>3\sigma$) within 2026--2027
\item \textbf{1 is already in tension} (CPL parameters, 2.4$\sigma$)
\end{itemize}

The single most important near-term test is \textbf{Euclid's scale-dependent $S_8$} (Prediction E1), which is a zero-parameter prediction with 3--5$\sigma$ discriminating power and data expected late 2026. The single most important cross-check is the \textbf{CMB lensing $\times$ galaxy cross-correlation} (Prediction X1), which probes the same $G_\mathrm{eff}(k)$ through an independent channel.

The most urgent theoretical need is a \textbf{full DFD Boltzmann code} to convert the approximate psi-screen mapping into exact predictions, particularly for the DESI CPL parameters where the current 2.4$\sigma$ tension may be partly an artifact of the approximate mapping.

\end{document}
